\section{Shock}

\subsection{Introduction}

\subsection{An Ideal Fluid}

\subsection{Definition of a Shock}

\subsection{Conservation Laws for a Single Shock}
\subsubsection{Laboratory Fixed Coordinate Frame}
The derivation that follows assumes a {\bf laboratory-fixed reference
frame}, as shown in Fig.~\ref{fig:piston_shock_cv}.\footnote{A shock-fixed 
reference frame derivation then is developed in a subsequent section.}
Consider a rigid tube containing a rigid piston material at ambient
density $\rho_0$ and ambient pressure $P_0$ at time $t=t_0$.  

\begin{figure}[htb]
  \begin{center}
  %\begin{overpic}[width=\textwidth, grid, tics=10]{fig/piston_shock_cv}
  \begin{overpic}[width=\textwidth]{fig/piston_shock_cv}
    \put(14, 29){$O$}
    \put(95, 31){$x$}
    \put(15, 38){$y$}
    \put(12, 27){$z$}
    \put(23,  9.5){$U_P \cdot \Delta t$}
    \put(49,  9.5){$(U_P - U_P) \cdot \Delta t$}
    \put(44,  3.5){$U_S \cdot \Delta t$}
    \put( 2, 47){\footnotesize piston at time $t=t_0$}
    \put(28, 47){\footnotesize piston at time $t=(t_0 + \Delta t)$}
    \put(58, 47){\footnotesize shock with speed $U_S$ at $t=(t_0 + \Delta t)$}
    \put(49, 30){$\rho_1, P_1, U_P$}
    \put(85, 30){$\rho_0, P_0$}
  \end{overpic}
  \end{center}
  \caption{Piston moving left-to-right along the $x$-axis, with 
  piston speed $U_P$, causing shock with shock speed $U_S$.}
  \label{fig:piston_shock_cv}
\end{figure}

Assume at time $t=t_0 + \Delta t$, the piston moves left-to-right
along the $x$-axis with speed $U_P$, measured in the fixed 
laboratory reference frame.  
Assume the piston speed $U_P$ causes a shock to propagate through
the material with speed $U_S$.  We require the following:
\begin{itemize}
  \item $0 < U_P < +\infty$
  \item $0 < U_S < +\infty$
  \item $U_P < U_S$
\end{itemize}  
From these considerations, we conclude that the material to the left of the
shock wave in Fig.~\ref{fig:piston_shock_cv} moves to the right at speed
$U_S$.  The material to the right of the shock wave remains at 
rest.\footnote{Furthermore, we assume the material throughout the tube and
for all time is inviscid, non-diffusive, and non-reacting.  We assume 
an adiabatic boundary, thus no
heat conduction occurs through the boundary of the tube walls or piston.
We assume no heat transfer from radiation occurs across the shock.  Finally, 
we assume body forces from gravity and electromagnetism are negligible.}

Since the piston compresses the material in the tube, we know
\begin{itemize}
\item $\rho_0 < \rho_1$.
\end{itemize}

Allow some time interval $\Delta t$ to elapse, as illustrated 
in Fig.~\ref{fig:piston_shock_cv}.  Consider the distances along the
$x$-axis, written as a multiplication of speeds and elapsed time.  
In the time interval $\Delta t$, the piston moves along the $x$-axis
a distance of $U_P \cdot \Delta t$.  
In the same time interval, the shock wave moves along the $x$-axis
a distance of $U_S \cdot \Delta t$.  Thus, the control volume at 
$t=(t_0 + \Delta t)$ has an axial length of $(U_S - U_P) \cdot \Delta t$.

From these axial lengths and piston cross-sectional area $A$, we can write
the following relationships:

\begin{itemize}
\item volume
\begin{itemize}
%
\item[] $V_0 = \left.V\right|_{t=t_0} = A \cdot U_S \cdot \Delta t$,
\item[] $V_1 = \left.V\right|_{t=(t_0 + \Delta t)} = A \cdot \left(
U_S - U_P\right) \cdot \Delta t$,
\end{itemize}
%
\item density
\begin{itemize}
\item[] $\rho_0 = \left.\rho\right|_{t=t_0}$,
\item[] $\rho_1 = \left.\rho\right|_{t=(t_0 + \Delta t)}$,
\end{itemize}
%
\item mass 
\begin{itemize}
\item[] $\rho_0 V_0 = \rho_0 \cdot A \cdot U_S \cdot \Delta t$, 
\item[] $\rho_1 V_1 = \rho_1 \cdot A \cdot \left(U_S - U_P\right) \cdot 
\Delta t$.
\end{itemize}
%
\end{itemize}

\subsubsection{Conservation of Mass}
Conservation of mass within the control volume requires that mass at
time $t=(t_0 + \Delta t)$ be unchanged from mass at time $t=t_0$.  
Thus
\be
 \rho_0 \cdot A \cdot U_S \cdot \Delta t =
 \rho_1 \cdot A \cdot \left(U_S - U_P\right) \cdot \Delta t, 
\ee
or
\be
 \boxed{\rho_0 U_S = \rho_1 \left(U_S - U_P\right).}
\ee

\subsubsection{Conservation of Momentum}
The driving force from the piston on the material is $(P_1 - P_0) \cdot A$.
Conservation of momentum requires $F = d(ma)/dt$.  Written over the time
interval $\Delta t$ and for constant mass, 
conservation of momentum can be written
\be
 F \cdot \Delta t = m \cdot \Delta U.
\ee
Rewriting this expression in terms of pressures and densities, we have
\be(P_1 - P_0) \cdot A \cdot \Delta t = \left(\rho_0 \cdot A \cdot U_S \cdot 
\Delta T \right) \; \left(U_P - 0\right) 
\ee
which gives
\be
\boxed{P_1 - P_0  = \rho_0 \; U_S \; U_P.}
\label{eq:cons_momentum}
\ee

\subsubsection{Conservation of Energy}
The work done on the control volume by the piston moving left-to-right
is given by
\be
W_P = \underset{\mbox{\footnotesize force}}{\underbrace{P_1 \cdot A}}
\cdot 
 \underset{\mbox{\footnotesize disp.}}{\underbrace{U_P \cdot \Delta t}}.
\ee
The kinetic energy increase of the control volume is simply the one-half times
the mass times the square of the change in velocity, thus
\be
KE = \half  \;
\underset{\mbox{\footnotesize mass}}{\underbrace{\left(\rho_0 \cdot A 
\cdot U_S \cdot \Delta t\right)}} \;
\underset{\mbox{\footnotesize velocity}}{\underbrace{\left(U_P - 0
\right)}}^2.
\ee
Finally, the increase in internal energy is just the internal energy 
{\bf per unit mass}, $e$, multiplied by the mass, 
\be
IE = (e_1 - e_0)  \;
\underset{\mbox{\footnotesize mass}}{\underbrace{\left(\rho_0
 \cdot A \cdot U_S \cdot \Delta t\right)}}.
\ee
Equating these quantities gives
\be
P_1 \cdot A \cdot U_P \cdot \Delta t = 
\half \; \rho_0 \cdot A \cdot U_S \cdot \Delta t \cdot U_P^2 + 
\left(e_1 - e_0\right) \; \rho_0 \cdot A \cdot U_S \cdot \Delta t
\ee
or simply
\be
\boxed{P_1 \; U_P = \rho_0 \; U_S \left(\frac{U_P^2}{2} + e_1 - e_0\right).}
\label{eq:cons_energy}
\ee
The foregoing equation can be written without kinematic variables 
$U_S$ and $U_P$. Define the specific volume as
\be
V \defe \frac{1}{\rho}.
\label{eq:specific_volume}
\ee
Next, note that conservation of mass can be rewritten as 
\be
\frac{U_P}{U_S} = 1 - \frac{\rho_0}{\rho_1}.
\label{eq:cons_mass_ratio}
\ee
Then, with (\ref{eq:cons_energy}),
\begin{align}
P_1 \; U_P 
&\overset{(\ref{eq:cons_momentum})}{=} 
\left(P_1 - P_0\right) \frac{U_P}{2} + \rho_0 \; U_S \; (e_1 - e_0), \\
&= -P_0 \; U_P + 2 \; \rho_0 \; U_S \; (e_1 - e_0), \\
(e_1 - e_0) &= \frac{U_P \; (P_1 + P_0)}{2 \; \rho_0 \; U_S}, \\
&\overset{(\ref{eq:cons_mass_ratio})}{=} 
\left(\frac{P_1 + P_0}{2\;\rho_0}\right) \; 
\left(1 - \frac{\rho_0}{\rho_1}\right), \\
&\overset{(\ref{eq:specific_volume})}{=} 
\left(\frac{P_1 + P_0}{2}\right) V_0 \left(1 - \frac{V_1}{V_0}\right), 
\end{align}
which simplifies to
\be
\boxed{e_1 - e_0 = \left(\frac{P_1 + P_1}{2}\right) \; (V_0 - V_1).}
\ee

\subsubsection{Alternative Formulations}
Now we develop the {shock-fixed reference frame} equations for mass, 
momentum, and energy conservation.

